\FloatBarrier
\section{Seeing the generated journeys and their GAF feedback}
\label{sec:g36-visual}

The delivered object is a timestamped coordinate trajectory. To make that output visible, Figures~\ref{fig:g36-journey-bus}--\ref{fig:g36-gaf-taxi_or_uber} connect the expanded model's generated journeys to the exact two-view images used for feedback. These are completed the preliminary candidate-selection study outputs, not schematic curves or newly sampled examples chosen to improve a figure. Before rendering, the display rule was fixed to the morning episode, expanded native model seed 42, expanded classifiers and the first lexicographic request in each class jointly supported by the reference, map and four policies. All five classes receive the same paired trajectory/image presentation in the main text. The cases explain the output and feedback mechanism; one fixed request per class is not an estimate of typical quality. Population-level comparisons below retain all classes and physical endpoints.

\subsection{Geometry and motion are different parts of a trajectory}
The two maps in each journey panel display the independent SUMO realization and the Uniform diffusion output, with road lanes in grey and elapsed trip fraction in colour. Start and end markers locate the complete journey. The map geometry is shared by design; it does not demonstrate that GAF has learned road connectivity. Timing along that geometry is the generated quantity. Plotting only the XY curves would conceal differences between candidates that traverse the same route at different speeds.

The lower panels therefore show physical elapsed time, accumulated progress and speed. Pale progress curves retain all eight candidates; the coloured curves identify the four recorded selectors. This separates the diversity available before feedback from the single output delivered by each selector. The reference is an independently simulated realization for the request, not a future trajectory supplied to diffusion or a unique deterministic target that every synthetic sample should reproduce.

\fig{g36-journey-bus.pdf}{1}{Actual complete bus journeys under a fixed road request. Top: independent SUMO reference and Uniform diffusion output, with a shared local coordinate frame and within-trip time colouring; white circle and cross mark start and end. Bottom: all eight candidate progress curves (pale) and the four selected outputs, alongside reference motion (dashed). Speed averages complete one-second intervals. Different reference and generated durations remain visible; the four generated selections share one duration.}{g36-journey-bus}
The displayed bus also reveals a limitation that a map overlay would hide: its independent reference lasts 2,315 seconds, whereas the shared generated duration is 1,235.64 seconds. All selectors inherit that value, so none can correct the overall travel-time mismatch within this pool. Spatial selection changes the pattern of movement within the allotted duration. That is a different perturbation from learning a more realistic duration prior.

\subsection{The GAF images associated with those outputs}
Figure~\ref{fig:g36-gaf-bus} shows the actual images obtained from the same five trajectories. Each spatial RGB image consists of the inherited GASF/GADF/MTF channels; the velocity image contains the corresponding speed-derived fields. The generated spatial pixel hashes match the archived images used for scoring. Velocity panels use the exact classifier-input transform with its final ImageNet normalization undone for display. No image receives an individually optimized colour range.

\fig{g36-gaf-bus.pdf}{1}{Bus trajectories and their exact GAF representations. Rows are the SUMO reference and the recorded Uniform, spatial (S), velocity (V) and joint (SV) selections. The right column ties every image pair back to its generated motion. $p_S$ and $p_V$ are the stored target-class probabilities of the two classifiers; they are auxiliary scores, not independent realism measures. Repeated candidate choices intentionally produce repeated image rows.}{g36-gaf-bus}
In this bus pool, Uniform, V and SV retain the same candidate, while S retains another. The spatial target probability rises from 0.9915 to 0.9992 under S; the velocity probabilities remain very small for these displayed candidates. These observations are a concrete warning against interpreting high spatial class confidence as sufficient evidence of realistic whole-trip timing. They also explain why apparently repeated image rows are informative rather than a plotting error.

The motorcycle pair follows the same display contract as the bus. Figures~\ref{fig:g36-journey-motorcycle} and \ref{fig:g36-gaf-motorcycle} apply the identical case rule, transforms and display settings to a motorcycle. Comparing image patterns is useful for understanding what changes under selection, but visual texture alone cannot explain participant semantics or establish a latent causal mechanism. Those claims require controlled quantitative evidence.

\fig{g36-journey-motorcycle.pdf}{1}{Complete motorcycle generation under the same display protocol as the bus. Map agreement reflects the route interface; progress and speed expose the generated timing and motion. All four recorded policies and all eight pre-selection progress curves are retained, regardless of whether they appear closer to the reference.}{g36-journey-motorcycle}
\fig{g36-gaf-motorcycle.pdf}{1}{The motorcycle's exact spatial and speed image inputs, aligned with the corresponding motion curves. Both image axes follow sequence position, not map coordinates. All channels and candidate choices are produced by the same frozen transformations and selectors as the aggregate experiment. A higher classifier score is not presented as a physical-quality guarantee.}{g36-gaf-motorcycle}
\subsection{Private car, bicycle and taxi/ride-hailing}
The remaining classes are part of the same synthesis task, not secondary examples. Private cars dominate the training corpus, bicycles supply a distinct low-speed domain, and taxi/ride-hailing retains a service label that need not correspond to a unique local movement pattern. All three therefore matter when interpreting what a participant-conditioned generator has learned and what its GAF classifiers can add.

The private-car example has very high target probabilities in both classifiers, yet the reference duration is 747.25 seconds and the generated duration is 834.07 seconds. The bicycle example shows a different representation mismatch: its Uniform spatial target probability is approximately $2.59\times10^{-8}$ while its velocity probability is saturated near one. Spatial selection changes the chosen image but still gives a low bicycle target score, approximately $3.15\times10^{-5}$. This is one generated-request diagnostic, not a contradiction of the classifier's development-set F1.

For the taxi/ride-hailing example, Uniform has spatial/velocity probabilities of approximately 0.307/1.000. Spatial selection raises its own probability near one but chooses a candidate whose velocity probability is only 0.082. Velocity and joint selection choose other outputs. This disagreement makes the proposed two-view mechanism concrete; it does not by itself establish that either view or their combination selects more realistic motion. The five-class physical evaluation is needed to answer that question.

\fig{g36-journey-private_car.pdf}{1}{Private car: complete independent reference and generated journeys, all eight progress curves and four recorded selections. Same fixed-case rule as the main visual results.}{g36-journey-private_car}
\fig{g36-gaf-private_car.pdf}{1}{Private car: exact spatial and speed GAF images of the corresponding reference and selected outputs, with motion and auxiliary classifier probabilities. Repeated selections are retained.}{g36-gaf-private_car}
\FloatBarrier

\fig{g36-journey-bycicle.pdf}{1}{Bicycle: complete independent reference and generated journeys, all eight progress curves and four recorded selections. Same fixed-case rule as the main visual results.}{g36-journey-bycicle}
\fig{g36-gaf-bycicle.pdf}{1}{Bicycle: exact spatial and speed GAF images of the corresponding reference and selected outputs, with motion and auxiliary classifier probabilities. Repeated selections are retained.}{g36-gaf-bycicle}
\FloatBarrier

\fig{g36-journey-taxi_or_uber.pdf}{1}{Taxi/ride-hailing: complete independent reference and generated journeys, all eight progress curves and four recorded selections. Same fixed-case rule as the main visual results.}{g36-journey-taxi_or_uber}
\fig{g36-gaf-taxi_or_uber.pdf}{1}{Taxi/ride-hailing: exact spatial and speed GAF images of the corresponding reference and selected outputs, with motion and auxiliary classifier probabilities. Repeated selections are retained.}{g36-gaf-taxi_or_uber}
\FloatBarrier

\subsection{What feedback can change in a fixed candidate pool}
Figure~\ref{fig:g36-eight-candidates} opens the bus pool before selection: all eight generated motions and both image views are shown. Policy tags identify which existing candidate each strategy retained. Terminal feedback cannot create a ninth motion profile or repair a flaw shared by every candidate. It can only alter which available realization is delivered. This is why the following sections evaluate physical properties of the selected outputs and compare them against Uniform under the same candidate budget.

\fig{g36-eight-candidates.pdf}{1}{The complete eight-candidate bus pool, with one-second speed and both actual image views for every candidate. Tags mark recorded selections; multiple tags mean the same output was retained by multiple policies. All candidates have the same requested route, class and duration. This panel exposes the available alternatives rather than displaying only the winning GAF image.}{g36-eight-candidates}

The new visual material changes the presentation, not the experimental outcomes. It adds no model fit, candidate search, teacher scoring or test access. The full scene/seed and class-level results below remain the basis for conclusions about generation quality and feedback utility.

\FloatBarrier
\section{Do the two GAF views recognize all five classes?}
\label{sec:expanded-recognition}
The current classifiers are trained on 29,895 complete native-demand legs and evaluated on a shared development population of 7,230 legs. The question here is whether each representation supplies usable class information for the synthesis task. Accuracy is 86.17\% for the spatial classifier and 89.13\% for the velocity classifier; macro F1 is 76.47\% and 90.34\%, respectively. Table~\ref{tab:g35-five-class-readers} reports every class rather than relying on the more numerous private cars.

\tab{Current GAF readers on the same development population. F1 values are percentages. Support counts are legs, not independent simulator scenes.}{g35-five-class-readers}{\begin{tabular}{lrrr}\toprule
Class & Development legs & Spatial F1 & Velocity F1 \\
Bus & 155 & 95.65 & 100.00 \\
Motorcycle & 1591 & 73.57 & 78.90 \\
Private car & 4436 & 95.20 & 92.23 \\
Bicycle & 668 & 72.10 & 99.48 \\
Taxi / ride-hailing & 380 & 45.85 & 81.07 \\
Macro F1 & 7,230 & 76.47 & 90.34 \\
\bottomrule\end{tabular}}

\fig{g35-confusion.pdf}{1}{Expanded-classifier development confusion matrices, normalized within each true class. The common five-class labels allow direct comparison. Spatial taxi/ride-hailing recognition remains weak despite high overall accuracy; speed-based recognition also retains confusion among motorized classes.}{g35-confusion}

Both views recognize buses strongly in this development cohort. Velocity F1 is also high for bicycles, whereas spatial F1 is lower. For private cars, spatial F1 exceeds velocity F1. Motorcycles and taxi/ride-hailing remain less well separated, especially in the spatial view. These differences motivate keeping both representations in the generation comparison, but they do not identify which selector will improve a particular class's movement.

The classifiers are fixed 50-epoch self-trained networks, not supplied group CNNs. Shared road and service templates, minority-class support and confidence calibration limit the interpretation of their scores. The earlier and current corpus comparisons are retained in Appendix~\ref{sec:corpus-reader-supplement} as data and optimization diagnostics. They are not a scaling-law contribution and do not replace generation evaluation.

\FloatBarrier
\section{Independent five-class generation evaluation}
\label{sec:expanded-generation}

Both new reference simulations, all eight candidate pools and all 48 selected-set evaluations completed. Each scene has 100 predeclared requests and 97 map-supported conditions, with the three unsupported bus routes retained in its denominator. The intersections with accepted reference trajectories contain 95 requests in the morning scene and 94 in the later scene. The eight full pools contain 6,208 generated trajectories in total. They are repeated conditional realizations, not 6,208 independent road environments. The finite cloud queue, including reference simulation, generation, scoring and evaluation, took approximately 64.9 minutes.

\tab{Out-of-sample replication coverage. Requests are shared public templates across scenes; paired reference support is reported separately from geometric support.}{g36-coverage}{\begin{tabular}{lrrr}\toprule
Scene & Requested & Map supported & Quality intersection \\
\midrule
Morning &100&97&95\\
Later &100&97&94\\
\bottomrule\end{tabular}}

The principal comparison uses the current generator with Uniform, S, V and SV selection. These four policies share the same eight candidates, route, class and duration for each request. Reference motion is opened only for evaluation. Old/new generator contrasts are preserved in Appendix~\ref{sec:corpus-generator-supplement}; their scene-dependent effects are a supporting domain-coverage diagnostic rather than the thesis's central result.

\FloatBarrier
\section{Does GAF feedback improve the generated motion?}
\label{sec:expanded-feedback}
The central comparison is between Uniform and the three GAF policies on identical current-generator pools. No selector receives an extra candidate or a different route or allotted duration. This evaluates what the trained representation contributes after class-conditioned diffusion has already generated possible journeys. Auxiliary target probabilities are reported in the visual panels; physical quality below is computed independently of those classifiers.

The predeclared native-acceleration distribution comparison remains mixed. Expanded velocity feedback versus Uniform changes the four scene/seed errors by $+8.7$, $+0.9$, $-9.6$ and $-4.7\%$. Thus, the two-seed acceleration benefit in the earlier fixed-duration development diagnostic does not reproduce uniformly across the new scenes and expanded components. The replication provides evidence about the boundary of that finding rather than invalidating its original, narrower population.

\fig{g36-metric-audit.pdf}{1}{Complete metric audit for expanded-teacher feedback relative to Uniform. Entries are percentage changes in equally weighted scene/seed mean errors; blue is lower and red is higher. W1 and request-level MAE address different targets. Controls are retained. This aggregate view is post hoc, with the individual cells and original protocol preserved.}{g36-metric-audit}

\tab{W1 change relative to Uniform (percent), with number of improving scene/seed cells in parentheses. Post-hoc equal-cell summary; lower is better. No significance test is implied}{g36-w1}{\begin{tabular}{lrrr}\toprule
Endpoint & Spatial S & Speed V & Dual SV \\
Duration [control] & +0.0 (0/4) & +0.0 (0/4) & +0.0 (0/4) \\
Observed arc [control] & -1.2 (4/4) & +0.3 (2/4) & +0.0 (2/4) \\
Longest steady movement & -4.9 (3/4) & -5.4 (3/4) & -6.2 (3/4) \\
Four-second speed change & -2.6 (3/4) & -0.5 (2/4) & -4.0 (2/4) \\
Mean speed [control] & +0.0 (2/4) & -0.0 (2/4) & +0.0 (2/4) \\
Mean absolute acceleration & -3.8 (3/4) & +0.3 (1/4) & -3.9 (2/4) \\
Stopped fraction & +3.1 (1/4) & +1.4 (2/4) & +3.2 (2/4) \\
95th-percentile speed & +2.0 (2/4) & -1.3 (2/4) & -5.8 (2/4) \\
Longest stop & -4.6 (3/4) & -6.3 (3/4) & -4.7 (3/4) \\
Stop episodes/minute & +5.2 (2/4) & +10.5 (1/4) & +6.3 (1/4) \\
Native speed distribution & +4.2 (1/4) & +1.0 (2/4) & +1.1 (3/4) \\
Native acceleration distribution & -2.0 (2/4) & -0.8 (2/4) & +0.0 (1/4) \\
\bottomrule\end{tabular}}

\tab{MAE change relative to Uniform (percent), with number of improving scene/seed cells in parentheses. Post-hoc equal-cell summary; lower is better. No significance test is implied}{g36-mae}{\begin{tabular}{lrrr}\toprule
Endpoint & Spatial S & Speed V & Dual SV \\
Duration [control] & +0.0 (0/4) & +0.0 (0/4) & +0.0 (0/4) \\
Observed arc [control] & -1.3 (3/4) & +0.7 (1/4) & +0.1 (1/4) \\
Longest steady movement & -3.6 (3/4) & -1.3 (2/4) & -3.3 (3/4) \\
Four-second speed change & -1.4 (3/4) & +0.2 (2/4) & -3.2 (3/4) \\
Mean speed [control] & -0.0 (4/4) & -0.0 (4/4) & -0.0 (3/4) \\
Mean absolute acceleration & +2.6 (1/4) & -0.9 (2/4) & -0.5 (3/4) \\
Stopped fraction & +0.4 (1/4) & -2.0 (3/4) & -0.4 (2/4) \\
95th-percentile speed & -1.4 (3/4) & +2.6 (1/4) & +0.5 (1/4) \\
Longest stop & -7.6 (4/4) & -3.6 (3/4) & -7.0 (3/4) \\
Stop episodes/minute & +1.8 (2/4) & -0.8 (2/4) & -0.2 (2/4) \\
\bottomrule\end{tabular}}


Tables~\ref{tab:g36-w1} and \ref{tab:g36-mae} display all recorded scalar endpoints as well as the native distribution measures where defined. The mean errors are averaged before their ratio is taken; an average percentage change would answer a different weighting question. W1 is not removed when a paired error appears more favourable. Nor are nearly invariant travel time and mean speed counted as useful feedback wins. Small observed-arc differences arise from temporal coordinate sampling on a common route and do not establish improved route planning.

The combined view identifies movement-specific effects. All three selectors reduce the equal-cell mean longest-stop and longest-steady-movement distribution errors. At the same time, all three increase mean stopped-fraction and stop-frequency distribution errors. Velocity feedback lowers longest-stop W1 by 6.3\% but raises stop-frequency W1 by 10.5\%. Spatial feedback lowers native-acceleration W1 by 2.0\% while raising native-speed W1 by 4.2\%. The equal-weight dual policy does not consistently dominate either single view. These paired benefits and costs belong together in the main interpretation.

\FloatBarrier
\section{Five-class motion effects and representation tradeoffs}
\label{sec:exploratory-stopping}
A participant-conditioned synthesis system must account for every supported class. Figure~\ref{fig:g36-five-class-motion} therefore places all five classes, all three GAF policies and six motion-distribution endpoints together. Table~\ref{tab:g36-five-class-native} gives the underlying absolute speed and acceleration errors, so a relative improvement can be read against its baseline scale. The aggregate errors average the four scene/seed cells equally; counts of improving cells are descriptive, not independent statistical tests. This class-level reading is post hoc and is not a rule for selecting a different policy for each class.

\fig{g36-five-class-motion.pdf}{1}{Complete five-class motion comparison for current-generator GAF feedback versus Uniform. Entries give the percentage change in mean W1 error and the number of improving scene/seed cells. Negative values are better. The same colour scale is used throughout; all classes and six motion-distribution endpoints are retained. Four cells contain two simulator scenes and two model seeds, not four independent scenes.}{g36-five-class-motion}
\tab{Absolute native-distribution W1 errors of the current generator and classifiers, averaged equally over the four scene/seed cells. Speed uses m/s and acceleration uses m/s$^2$; lower is better. Units must not be pooled}{g36-five-class-native}{\begin{tabular}{lrrrr}\toprule
Class / endpoint & Uniform & S & V & SV \\
Bus / speed & 0.544 & 0.594 & 0.532 & 0.558 \\
Bus / accel. & 0.089 & 0.093 & 0.092 & 0.094 \\
Motorcycle / speed & 1.233 & 1.180 & 1.205 & 1.218 \\
Motorcycle / accel. & 0.318 & 0.286 & 0.329 & 0.308 \\
Private car / speed & 1.475 & 1.584 & 1.460 & 1.465 \\
Private car / accel. & 0.183 & 0.181 & 0.177 & 0.174 \\
Bicycle / speed & 0.535 & 0.548 & 0.564 & 0.541 \\
Bicycle / accel. & 0.198 & 0.199 & 0.178 & 0.207 \\
Taxi / ride-hailing / speed & 1.309 & 1.404 & 1.385 & 1.372 \\
Taxi / ride-hailing / accel. & 0.140 & 0.150 & 0.144 & 0.146 \\
\bottomrule\end{tabular}}


\subsection{Bus: stopping structure and complete-trip timing}
All three selectors reduce bus longest-stop distribution error in all four cells, with equal-cell reductions of 12.7\% for S, 22.5\% for V and 13.4\% for SV. Stop-frequency errors also decrease on average. Native acceleration errors, however, increase on average for all three. These findings concern different parts of motion and cannot be merged into a claim that the bus trajectory is uniformly better. The fixed visual example further shows that selection cannot correct a travel-time error shared by the pool. The earlier paired longest-stop analysis is retained in Appendix~\ref{sec:exploratory-stopping-supplement}; it does not establish correct station-stop location.

\subsection{Motorcycle: acceleration rather than a stopping-only narrative}
Spatial feedback reduces motorcycle native-acceleration W1 by 10.3\% in the equal-cell summary, with all four cells improving. Native-speed W1 decreases by 4.3\%, with three cells improving. Velocity feedback instead increases acceleration W1 by 3.4\% on average. These comparisons show why the two representations should be evaluated on motion quality even when the velocity classifier has higher development F1. Motorcycle longest-stop W1 increases under all three policies, retaining a concrete cost alongside the acceleration finding.

\subsection{Private car: common-class performance and stop-frequency cost}
Joint feedback reduces private-car acceleration W1 by 5.0\% and longest-stop W1 by 8.9\%, with three cells improving for each. Native-speed W1 is nearly unchanged on average under V and SV, while spatial selection raises it by 7.4\%. All three policies raise stop-frequency W1 by approximately 24\%. The numerous private-car training examples and strong classifier recognition therefore do not make every physical property easy to improve. This class is essential to the system's coverage and to the interpretation of its macro averages.

\subsection{Bicycle: a velocity-feedback gain with a speed cost}
Velocity feedback reduces bicycle native-acceleration W1 by 9.8\%, with all four cells improving, while it increases native-speed W1 by 5.4\%, with no cell improving. Spatial feedback reduces longest-stop W1 by 18.8\% in the equal-cell summary, but does not reproduce the velocity selector's acceleration benefit. This pattern is especially informative because a nearly perfect velocity development classifier does not produce a uniformly better synthetic bicycle trajectory. The task is to reproduce a within-class motion distribution, not merely a recognizable bicycle image.

\subsection{Taxi/ride-hailing: representation disagreement and transfer limits}
Taxi/ride-hailing is retained despite weaker spatial recognition and adverse generation effects. All three policies raise mean native-speed and acceleration W1. Spatial selection raises stop-frequency W1 by 45.9\%, with no cell improving; velocity and joint policies also increase its mean error. The fixed candidate example shows strong cross-view disagreement, but the aggregate physical outcomes do not establish a useful selector for this class. A service label may overlap with ordinary car motion, and the current evidence does not identify which geographical or temporal feature causes the classifier's ranking.

\subsection{What the five-class evidence supports}
The spatial-motorcycle acceleration and velocity-bicycle acceleration patterns broaden the mechanism question beyond bus stopping. They are specific observations to replicate with new simulator episodes, alongside the adverse endpoints and every class. They do not justify a post-hoc class-specific policy switch, establish complementary latent factors, or isolate GAF's advantage over a matched non-image class classifier. That last comparison remains necessary to attribute utility to the representation rather than generic class-based selection.

Full-pool pair diversity describes the variation available before selection. With one retained sample per request, conditional pair diversity is undefined, not zero; across-request dispersion cannot replace it. Geometric support remains 97 of 100 requests in each scene and belongs to the common route interface. The result is a five-class trajectory-synthesis study with representation-dependent motion effects, rather than a claim that more data alone or one feedback score solves generation quality.
